\documentclass[a4paper]{article}

\usepackage[fontset=fandol]{ctex}
\usepackage{amsmath, amssymb}
\usepackage{graphicx}
\usepackage[margin=2.45cm]{geometry}
\usepackage[most]{tcolorbox}
\usepackage{hyperref}
\usepackage{booktabs}
\usepackage{float}
\usepackage{array}
\usepackage{xcolor}
\usepackage{enumitem}
\usepackage{caption}

\setlength{\parindent}{2em}
\setlength{\parskip}{0.35em}
\setlength{\emergencystretch}{3em}
\linespread{1.06}
\sloppy
\raggedbottom
\vfuzz=4pt

\newcolumntype{L}[1]{>{\raggedright\arraybackslash}p{#1}}
\newcolumntype{C}[1]{>{\centering\arraybackslash}p{#1}}

\DeclareMathOperator*{\argmin}{arg\,min}
\DeclareMathOperator*{\argmax}{arg\,max}
\newcommand{\E}{\mathbb{E}}
\newcommand{\policy}{\pi_\theta}
\newcommand{\traj}{\tau}

\hypersetup{
    colorlinks=true,
    linkcolor=blue!60!black,
    urlcolor=blue!60!black
}

\setlist[itemize]{leftmargin=2.2em,itemsep=0.22em,topsep=0.25em}
\setlist[enumerate]{leftmargin=2.4em,itemsep=0.22em,topsep=0.25em}
\setlist[description]{leftmargin=3.0em,itemsep=0.18em,topsep=0.25em}

\newtcolorbox{knowledgebox}[1]{
    enhanced, colback=blue!5!white, colframe=blue!75!black, colbacktitle=blue!75!black,
    coltitle=white, fonttitle=\bfseries, title={#1},
    attach boxed title to top left={yshift=-2mm, xshift=2mm},
    boxrule=1pt, sharp corners
}

\newtcolorbox{importantbox}[1]{
    enhanced, colback=yellow!10!white, colframe=yellow!80!black, colbacktitle=yellow!80!black,
    coltitle=black, fonttitle=\bfseries, title={#1}, sharp corners
}

\newtcolorbox{warningbox}[1]{
    enhanced, colback=red!5!white, colframe=red!75!black, colbacktitle=red!75!black,
    coltitle=white, fonttitle=\bfseries, title={#1}, sharp corners
}

\newcommand{\notetitle}{强化学习概述（二）：Policy Gradient}
\newcommand{\notesubtitle}{Return Estimation, Baseline, On-policy, Off-policy, PPO, and Exploration}
\newcommand{\noteauthors}{基于李宏毅老师《机器学习 2025》课程内容整理}
\newcommand{\notedate}{\today}
\newcommand{\videochannel}{李宏毅深度学习课堂（Bilibili）}
\newcommand{\videoduration}{约 41 分 13 秒}
\newcommand{\videourl}{https://www.bilibili.com/video/BV1TAtwzTE1S?p=47}

\newcommand{\framefig}[4]{%
\begin{figure}[H]
    \centering
    \includegraphics[width=#1\textwidth]{#2}
    \caption{#3\protect\footnotemark}
\end{figure}
\footnotetext{视频画面时间区间：#4。}
}

\begin{document}
\pagestyle{empty}

\begin{center}
    \vspace*{1.0cm}
    {\Huge\bfseries \notetitle\par}
    \vspace{0.35cm}
    {\LARGE \notesubtitle\par}
    \vspace{0.75cm}
    \includegraphics[width=0.62\textwidth]{video.jpg}\par
    \vspace{0.75cm}
    {\large \noteauthors\par}
    {\large \notedate\par}
    \vspace{0.75cm}
    \begin{tcolorbox}[width=0.86\textwidth,center,sharp corners,
                      colback=black!3!white,colframe=black!50!white]
        \centering
        \textbf{视频信息}\\[3pt]
        频道：\videochannel \\
        时长：\videoduration \\
        链接：\href{\videourl}{\videourl}
    \end{tcolorbox}
\end{center}

\newpage
\tableofcontents
\newpage
\pagestyle{plain}

% =====================================================================
\section{从 P46 留下的问题开始}
% =====================================================================

上一讲把 policy gradient 的直觉写成加权 cross-entropy：如果在 state $s_n$ 下采取 action $a_n$ 是好事，就给它正权重；如果是坏事，就给它负权重。训练 actor 的目标可以写成
\[
L(\theta)=\sum_{n=1}^{N}A_n e_n
=-\sum_{n=1}^{N}A_n\log \pi_\theta(a_n\mid s_n),
\]
其中 $e_n=-\log\pi_\theta(a_n\mid s_n)$。问题是：这些 $(s_n,a_n)$ 从哪里来？权重 $A_n$ 又怎么知道？

\framefig{0.92}{figures/fig01_pg_open_questions_from_p46.jpg}
{P46 结尾留下两个问号：训练 actor 的 state-action pair 和每个 pair 的权重 $A_n$ 从哪里来？}
{00:00:00--00:00:15}

P47 就是在回答这两个问题。最自然的答案是：让 actor 自己去和环境互动，收集许多 episode。每一次互动都会产生 state、action 和 reward。state-action pair 直接来自 trajectory；而 $A_n$ 则要从 reward 中估计。接下来课程用 Version 0、1、2、3 逐步改进 $A_n$ 的定义。

\begin{importantbox}{policy gradient 的核心任务}
    policy gradient 要把“整条 trajectory 的好坏”转换成“每个动作该被鼓励还是抑制”的权重。这个权重就是本节反复出现的 $A_n$，它可以理解为某个动作相对于平均水平的好坏程度。
\end{importantbox}

\subsection{本章小结}

P47 的主线不是重新定义 actor，而是补上 P46 中缺失的训练信号来源：用当前 actor 收集 trajectory，再从 reward 中构造 $A_n$，最后用加权 cross-entropy 更新 actor。

% =====================================================================
\section{Version 0：即时 reward 当作动作评价}
% =====================================================================

最简单的做法是收集许多 episode，把每一步的 state-action pair 作为训练资料：
\[
\mathcal{D}=\{(s_1,a_1),(s_2,a_2),\ldots,(s_N,a_N)\}.
\]
然后直接把该步得到的即时 reward 当作权重：
\[
A_t=r_t.
\]

\framefig{0.92}{figures/fig02_version0_collect_state_action_pairs.jpg}
{Version 0：actor 与环境互动，收集许多 $(s_t,a_t)$ pair 作为训练资料。}
{00:04:15--00:04:30}

\framefig{0.92}{figures/fig03_version0_immediate_reward_short_sighted.jpg}
{Version 0 把每一步的 $A_t$ 直接设为即时 reward $r_t$，因此是一个 short-sighted version。}
{00:07:15--00:07:30}

这个版本很直观：某一步得到正 reward，就鼓励这一步动作；某一步得到负 reward，就抑制这一步动作。但它的问题也很明显：动作的好坏不一定在同一时刻显现。很多动作是为未来铺路，短期看没有 reward，长期却很重要。

\subsection{reward delay 与短视问题}

在 RL 中，一个 action 会影响之后的 observation，也会影响之后能得到的 reward。若只看 $r_t$，actor 会变得短视。课程用 Space Invader 说明：如果只有“开火击中外星人”那一步有正 reward，而向右移动暂时没有分数，那么 Version 0 会学到“只有开火是好动作”，却学不到先移动再开火的策略。

\framefig{0.92}{figures/fig04_reward_delay_problem.jpg}
{reward delay：动作会影响后续观察与后续 reward，actor 有时必须牺牲即时 reward 来换取长期 reward。}
{00:08:30--00:08:45}

\framefig{0.92}{figures/fig05_space_invader_short_sighted_fire.jpg}
{Space Invader 中，若只把即时 reward 当评价，actor 可能只学会不断开火，而不会学习先移动到合适位置。}
{00:09:45--00:10:00}

\begin{warningbox}{即时 reward 不是动作价值}
    $r_t$ 只描述当前动作后的即时反馈，不等于 $a_t$ 对整局结果的贡献。若任务存在延迟回报，只用 $r_t$ 会造成错误 credit assignment。
\end{warningbox}

\subsection{本章小结}

Version 0 的训练资料来源很自然，但 $A_t=r_t$ 太短视。它忽略了动作对未来状态和未来 reward 的影响，因此无法处理 reward delay。下一步要把未来 reward 也纳入动作评价。

% =====================================================================
\section{Version 1 与 Version 2：累计与折扣累计 reward}
% =====================================================================

Version 1 把一个动作之后所有 future reward 加起来，作为该动作的评价。对第 $t$ 步动作，定义 return-to-go：
\[
G_t=\sum_{n=t}^{N}r_n.
\]
然后令
\[
A_t=G_t.
\]
这比即时 reward 合理：若 $a_t$ 暂时没有 reward，却让未来拿到高分，它仍然会得到正评价。

\framefig{0.92}{figures/fig06_version1_cumulative_reward.jpg}
{Version 1：用从当前时间点到 episode 结束的累计 reward $G_t$ 评价动作。}
{00:12:15--00:12:30}

不过，Version 1 又有另一个问题：很遥远的 reward 也被完全归功于早期动作。若 episode 很长，把 $r_N$ 全额算到 $a_1$ 头上，未必合理。早期动作当然可能影响未来，但影响通常会随着时间变远而减弱。

\subsection{折扣因子}

Version 2 引入 discount factor $\gamma<1$，让更远的 reward 权重更小：
\[
G_t'=\sum_{n=t}^{N}\gamma^{n-t}r_n
=r_t+\gamma r_{t+1}+\gamma^2 r_{t+2}+\cdots.
\]
然后令
\[
A_t=G_t'.
\]

\framefig{0.92}{figures/fig07_version2_discounted_return.jpg}
{Version 2：用折扣因子 $\gamma<1$ 降低远期 reward 的权重，避免把太遥远的结果完全归功于当前动作。}
{00:14:45--00:15:00}

折扣因子的作用有两层。第一，它符合“越近的结果越可能与当前动作有关”的直觉。第二，它让数值更稳定，尤其在 episode 很长甚至无限长时，折扣和可以保持有限。常见 RL 文献中，$\gamma$ 也表达 agent 对未来的重视程度。

\begin{knowledgebox}{return、reward、advantage 的关系}
    reward 是单步反馈，return 是从某步开始累积的未来反馈，advantage 则通常表示某个动作比基准预期好多少。本节的 $A_t$ 先从 reward、return 开始，最后通过减 baseline 接近 advantage 的形式。
\end{knowledgebox}

\subsection{本章小结}

Version 1 用未来累计 reward 解决短视问题，但可能过度归功于远期结果。Version 2 加入折扣因子，让未来 reward 的贡献随时间衰减。这已经比即时 reward 合理得多，但仍然缺少“相对好坏”的概念。

% =====================================================================
\section{Version 3：减 baseline，得到相对好坏}
% =====================================================================

课程接着指出：reward 的好坏是相对的。如果一次考试你拿 60 分，但全班平均 50 分，那 60 分不错；如果全班平均 80 分，那 60 分就很差。RL 中也是一样，若某个游戏里每个动作都得到正分，只看正 reward 会让所有动作都被鼓励，actor 不知道哪些动作其实只是“低于平均”。

Version 3 因此引入 baseline $b$：
\[
A_t=G_t'-b.
\]
这样 $A_t$ 可以有正有负。若 $G_t'$ 高于 baseline，动作被鼓励；若低于 baseline，动作被抑制。

\framefig{0.92}{figures/fig08_version3_baseline_relative_reward.jpg}
{Version 3：从折扣累计 reward 中减去 baseline $b$，让 $A_t$ 同时有正值和负值。}
{00:22:15--00:22:30}

在标准 policy gradient 中，这个 $A_t$ 往往被称为 advantage。最理想的形式是
\[
A^\pi(s_t,a_t)=Q^\pi(s_t,a_t)-V^\pi(s_t),
\]
其中 $Q^\pi(s_t,a_t)$ 表示在 $s_t$ 做 $a_t$ 后的预期 return，$V^\pi(s_t)$ 表示在 $s_t$ 按当前 policy 正常行动的平均预期 return。课程在这里先不展开 $Q$ 和 $V$，而是用 baseline 的直觉说明：动作好不好要看它是否比通常情况更好。

\begin{importantbox}{baseline 的两个作用}
    baseline 不改变“相对鼓励哪个动作”的核心方向，却能让权重有正有负，并降低梯度估计方差。直觉上，它把“拿到多少分”变成“比预期多拿多少分”。
\end{importantbox}

\subsection{本章小结}

Version 3 把 $A_t$ 从绝对 return 改成相对 return。减 baseline 后，actor 不再盲目鼓励所有正 reward 动作，而是鼓励高于基准的动作、抑制低于基准的动作。这是 policy gradient 稳定训练的重要思想。

% =====================================================================
\section{Policy Gradient 的训练流程}
% =====================================================================

把前面的 $A_t$ 估计放进训练流程，就得到本节的 policy gradient 框架。先初始化 actor 参数 $\theta^0$。在第 $i$ 次训练迭代中，用当前 actor $\theta^{i-1}$ 与环境互动，收集 trajectory；计算每个 state-action pair 的 $A_n$；构造 loss：
\[
L(\theta)=\sum_{n=1}^{N}A_n e_n
=-\sum_{n=1}^{N}A_n\log\pi_\theta(a_n\mid s_n);
\]
再做一次梯度更新：
\[
\theta^i \leftarrow \theta^{i-1}-\eta\nabla_\theta L(\theta).
\]

\framefig{0.92}{figures/fig09_policy_gradient_algorithm.jpg}
{policy gradient 流程：用 actor 与环境互动收集资料，计算 $A_n$，再用加权 loss 更新参数。}
{00:23:30--00:23:45}

\subsection{为什么只能更新一次}

课程特别强调一个与普通 supervised learning 不同的地方：在 supervised learning 中，训练资料通常先收集好，然后可以反复训练很多 epoch；但在 on-policy policy gradient 中，用 $\theta^{i-1}$ 收集到的数据，原则上只适合用来更新 $\theta^{i-1}$ 一次。更新成 $\theta^i$ 后，新的 actor 会产生不同的 trajectory，旧资料可能不再代表新 actor 的行为分布。

\framefig{0.92}{figures/fig10_update_once_recollect_data.jpg}
{on-policy PG 中，每次更新参数后通常需要重新收集整批训练资料；旧资料只用于一次更新。}
{00:26:00--00:26:15}

这就是 RL 训练很花时间的一个原因。每次参数更新都要与环境互动、重新采样 trajectory、重新计算 $A_n$。如果环境是真实机器人或真实用户系统，收集资料还可能非常昂贵。

\subsection{旧 trajectory 为什么不能随便复用}

假设旧 actor $\theta^{i-1}$ 在前几步和新 actor $\theta^i$ 做了同样动作，后面仍然可能走向完全不同的 trajectory。旧资料中某个动作后续得到高 reward，不代表新 actor 采取同一前缀动作后也会看到相同后续状态。因为 policy 改了，未来动作分布也改了。

\framefig{0.92}{figures/fig11_new_policy_may_not_observe_old_tail.jpg}
{旧 actor 收集到的 trajectory 后半段，新的 actor 不一定会观察到；这就是旧数据难以直接复用的原因。}
{00:29:45--00:30:00}

\begin{warningbox}{on-policy 的样本效率问题}
    on-policy 方法通常更直接、更符合推导假设，但样本效率较低。因为每次 policy 更新后，旧数据就变得不完全匹配，需要重新与环境互动采样。
\end{warningbox}

\subsection{本章小结}

policy gradient 的训练流程可以看成“采样、评估、更新、再采样”。它与监督学习最大的流程差异是数据收集在训练循环内部，而不是一次性固定在循环外部。这也引出 on-policy 与 off-policy 的区别。

% =====================================================================
\section{On-policy、Off-policy 与 PPO}
% =====================================================================

on-policy 指的是：要训练的 actor 和用来与环境互动、收集资料的 actor 是同一个。off-policy 则允许两者不同：一个 actor 负责收集经验，另一个 actor 利用这些经验学习。off-policy 的好处很直接：可以复用旧资料，不必每次更新都重新收集完整 trajectory。

\framefig{0.92}{figures/fig12_on_policy_vs_off_policy.jpg}
{on-policy 与 off-policy 的定义：训练用 actor 和互动收集资料的 actor 是否相同。}
{00:31:00--00:31:15}

\framefig{0.92}{figures/fig13_off_policy_reuse_old_trajectory.jpg}
{off-policy 试图用旧 actor 收集的 trajectory 来训练新 actor，以减少重复采样成本。}
{00:32:15--00:32:30}

但 off-policy 也有困难：若训练 actor 与互动 actor 差太多，旧经验对新 actor 的意义会打折扣。某个动作对旧 actor 是好经验，不代表新 actor 在相同状态下会延续到同样未来。严格算法需要处理 distribution mismatch，例如 importance sampling、trust region 或 clipped objective 等。

\subsection{PPO 的位置}

课程提到 Proximal Policy Optimization，简称 PPO，是今天常用的 off-policy / 近似 off-policy 思路之一。PPO 的核心精神是：训练 actor 要知道自己和收集资料的 actor 差多少，不要一次更新得太远。若新旧 policy 差异过大，旧数据就不可靠。

\framefig{0.92}{figures/fig14_ppo_actor_difference.jpg}
{PPO 的直觉：要训练的 actor 必须知道自己与收集资料的 actor 有多不同，并限制更新幅度。}
{00:33:30--00:33:45}

虽然课程没有展开 PPO 推导，但可以用一行公式理解其 clipped objective 的味道。设旧 policy 为 $\pi_{\theta_{\mathrm{old}}}$，新 policy 为 $\pi_\theta$，概率比为
\[
\rho_t(\theta)=
\frac{\pi_\theta(a_t\mid s_t)}
{\pi_{\theta_{\mathrm{old}}}(a_t\mid s_t)}.
\]
PPO 会避免 $\rho_t$ 偏离 $1$ 太远，防止新 policy 对旧 trajectory 做过度解释。

\begin{knowledgebox}{为什么 PPO 常见}
    PPO 通常比最朴素 policy gradient 稳定，又比更复杂的 trust-region 方法实现简单，因此在许多连续控制和模拟环境中成为常用 baseline。课程后面的 DeepMind 与 OpenAI 例子都与 PPO 相关。
\end{knowledgebox}

\subsection{本章小结}

on-policy 方法数据匹配但耗样本；off-policy 方法想复用旧数据，但必须处理旧 actor 与新 actor 的差异。PPO 的核心动机就是在可复用数据和不过度偏离旧 policy 之间取得平衡。

% =====================================================================
\section{Exploration：为什么 actor 要有随机性}
% =====================================================================

课程最后回到一个实践问题：数据从哪里来？若 actor 只会执行当前认为最好的动作，它可能永远不会尝试其他动作，也就永远不知道其他动作是否更好。例如 Space Invader 中，如果 actor 从不尝试开火，就不会知道开火能带来 reward。为了收集有价值的训练资料，actor 在 data collection 时必须有随机性。

\framefig{0.92}{figures/fig15_exploration_need_randomness.jpg}
{exploration：actor 在收集资料时需要随机性；这也是 policy 输出动作分布、从中采样的重要理由。}
{00:36:00--00:36:15}

常见 exploration 方法包括：
\begin{itemize}
    \item 按 $\pi_\theta(a\mid s)$ 采样动作，而不是总取最大概率动作；
    \item 提高输出分布的 entropy，让动作概率更分散；
    \item 在参数或动作上加入噪声；
    \item 在值函数方法中使用 $\epsilon$-greedy 等探索策略。
\end{itemize}

本节以 policy gradient 为主，所以重点是：采样动作不是实现细节，而是 RL 学习机制的一部分。没有探索，训练资料会太窄；资料太窄，actor 就很难发现更好的策略。

\subsection{PPO 在控制任务中的例子}

课程展示了 DeepMind 与 OpenAI 的 PPO 控制任务视频。例如让虚拟人学习走路、跑步、跨越障碍或在球场环境中运动。它们看起来有时很滑稽，但核心是同一个 RL 问题：actor 通过与模拟环境互动，逐渐学会能带来更高 reward 的动作序列。

\framefig{0.92}{figures/fig16_deepmind_ppo_locomotion.jpg}
{DeepMind PPO 控制任务示例：虚拟角色通过 reward 学习行走与跨越地形。}
{00:38:30--00:38:45}

\framefig{0.92}{figures/fig17_openai_ppo_locomotion.jpg}
{OpenAI PPO 控制任务示例：在模拟环境中训练连续控制策略，展示 policy gradient 方法的应用场景。}
{00:39:45--00:40:00}

这些例子也呼应了本讲标题中的“修课心情”：RL 训练常常充满不稳定、反复尝试、结果看起来怪但又逐渐变好。算法层面是 exploration 与 policy update，实践层面则是大量调参、重跑和观察曲线。

\subsection{走向 Actor-Critic}

P47 最后回到大纲，下一部分是 Actor-Critic。它会继续处理一个核心问题：如何更好地估计 $A_t$？单纯用 sampled return 方差很大，若能学习一个 critic 来估计 value 或 advantage，就可以让 actor 的更新更稳定。

\framefig{0.92}{figures/fig18_next_actor_critic_outline.jpg}
{P47 结束后，大纲进入 Actor-Critic；它会用 critic 来改进 actor 更新所需的价值估计。}
{00:40:45--00:41:00}

\subsection{本章小结}

exploration 让 actor 有机会发现未知动作的价值；PPO 等方法则在 policy 更新稳定性和资料复用之间折中。下一步 Actor-Critic 会把“如何估计 $A_t$”这个问题继续推进，用 critic 降低 policy gradient 的方差。

% =====================================================================
\section{总结与延伸}
% =====================================================================

这一讲把 P46 的直觉补成了可执行的 policy gradient 流程。首先，让 actor 与环境互动，收集 trajectory；其次，从 reward 中为每个 state-action pair 构造权重 $A_t$；最后，用
\[
L(\theta)=-\sum_t A_t\log\pi_\theta(a_t\mid s_t)
\]
更新 actor。核心变化在于 $A_t$ 的定义：Version 0 用即时 reward，太短视；Version 1 用未来累计 reward，解决 reward delay；Version 2 加入 discount factor，降低远期 reward 的影响；Version 3 减 baseline，把绝对分数变成相对好坏。

整节课最重要的概念压缩如下：
\begin{description}
    \item[return-to-go] $G_t=\sum_{n=t}^{N}r_n$，用当前之后的总 reward 评估动作。
    \item[discounted return] $G_t'=\sum_{n=t}^{N}\gamma^{n-t}r_n$，让远期 reward 衰减。
    \item[baseline / advantage] $A_t=G_t'-b$，鼓励高于基准的动作，抑制低于基准的动作。
    \item[on-policy] 用当前 actor 收集资料并训练当前 actor，数据匹配但耗样本。
    \item[off-policy / PPO] 尝试复用旧资料，但要控制新旧 actor 的差异。
    \item[exploration] actor 必须随机尝试动作，否则无法知道未尝试动作是否更好。
\end{description}

\begin{importantbox}{从 P46 到 P47 的完整闭环}
    P46 告诉我们可以用加权 cross-entropy 控制 actor；P47 告诉我们权重从哪里来：它来自 trajectory 中的 future reward，并通过 discount 与 baseline 修正。这样，长期 reward 被转换成了局部动作概率的更新方向。
\end{importantbox}

后续 Actor-Critic 会继续修正 policy gradient 的高方差问题。直觉上，actor 负责行动，critic 负责评价；critic 若能估计某个 state 或 action 的价值，就能给 actor 更可靠的 $A_t$。这也是强化学习从“靠采样回报更新”走向更稳定算法的重要一步。

\end{document}
